\section{Returning injectivity to the uncut Hecke point}

% BEGIN M87U1
\paragraph{M87U1. The string is present on every vector over the good point.}
Fix an original localization $i$, a good upper point $x$, and a nonzero
$f\in E_x$. Use the same row $j+1=j(i)$, tuple $u_0$ mapping to
$\mu(i)$, and upper-model inclusion $A$ of L86M1. The relation
$\mu(i)=\mu_{\rm At}(w(i))$ puts this tuple in the used range.
The vector $y=Af$ is nonzero by the norm identity of L86M1.
In the function model it is characterized uniquely by
$\operatorname{ext}_V(y)=X.\mathrm{fn}(f)$.

Choose a witnessing lower list $\Sigma$ and upper label $P_0$ for $x$.
The decomposition of $E_x$ is by all upper labels $P$ above $x$.
By H86P1 and H86P3, their base-change lists agree up to permutation
with the augmentation of $\Sigma$ for $P_0$. At every good split
place all these labels therefore give the same unramified parameter.
At an inert place every upper label is spherical. Its unitary
Satake parameter is the relative Weyl orbit of the torus tuple.
The base-change multiset there consists of the entries and their
inverses, with the fixed central entry in odd rank. Removing that
entry, when present, leaves the relative Weyl orbit uniquely
determined: the Weyl group permutes the inverse pairs and exchanges
their members. The multiplicities of self-inverse entries determine
the number of self-inverse pairs. Thus all labels also give the
same spherical character at each such place, with the splitting
normalization of the common group and frame.

For each allowed Hecke label take its eigenvalue from $P_0$.
The original finite representation on a spherical summand acts
through its Satake character on hyperspecial invariants; hence the
operator acts by that eigenvalue on the whole sum $E_x$.
The finite coset formula intertwines this action with the operator
on $y$: extension takes a right translate to the corresponding
translate of the extended form, and all character corrections
on the hyperspecial subgroup are one. Changing the representatives
therefore preserves the sum. This applies to any tuple of operators
representing the indicated Hecke labels. Assign arbitrary values
to the unused labels. The resulting system $s$ makes $y$ a
simultaneous eigenvector. At every good place take $B_w$ to be
the unramified base-change multiset just computed. The split
Satake formula and the inert relative Satake formula give the
defining eigenvalues of this same $s$, with no choice of an
irreducible summand containing the vector $f$.

The good-point list has a single cuspidal lower block by R86G4;
it augments by the specified length-$t$ character block. The
base-change comparison H86P1 is at all local places in its
unramified clauses. In particular, at the model's good places
the spherical rank-$r+t$ sum gives unramified constituents in
each summand (inertia and the monodromy operator restrict to
each direct summand). The remainder has $r$ roots of absolute
value at most $q_w^{(r-1)/2}$ by S86A2, and the appended
entries are $\chi_V(\varpi_w)q_w^{(t-1)/2-k}$ for $0\le k<t$.
Here the agreement of characters on the used tuple and common
coefficient embedding is M84T1 and T86C3. A place is good exactly
when neither it nor its conjugate is in the level and twist
support, so no ramified place has been inserted into this
calculation. This is the string property of D87C3 for the
same $y$. It holds before using a downward nonvanishing assertion.
% END M87U1

% BEGIN M87U2
\paragraph{M87U2. The same vector is fixed by the actual good idempotent.}
At the designated split places take the lower finite parameter
from the list $\Sigma$ of M87U1. It defines a parameter of
the member's type list: its archimedean weight has the prescribed
signature, its rank and its split localizations are those
of the lower list, and its designated labels satisfy the good
root tests by the choice of $x$. These are the signature, parameter
and local-label identifications in S86A1--4 and U86T3--4.
Every upper label above $x$ realizes the same augmentation:
permutation of the global list preserves all its local sums.
The native description of $M_t$ by its parameter types
therefore puts $E_x$ inside the parameter summand for this
single good parameter.

To see the projector on that summand, decompose its finite-group
action into irreducibles. At a type prime the parameter comparison
of I86A3 and the simultaneous reading T86C4 specify the
finite Green series of every occurring irreducible, equivalently
its polynomial label. The tuple of these series satisfies the
good tests. The good central idempotent is the sum of the products
of their series idempotents over all allowed tuples. Orthogonality
makes it the identity on each constituent in question and hence
on the entire parameter summand. Under the realization of $A$
each finite group element is a right translation corrected by the
level character. Changing its parahoric lift changes the translation
by a level element and the character by the inverse correction.
The selected coefficient embedding carries the coefficients of
the good idempotent to those defining the model idempotent.
Applying a scalar functional commutes with this finite sum.
Thus every $\ell(\operatorname{ext}_V(y))$ is fixed by the
model good idempotent. This is the condition $\mathrm{GoodVt}$
on the very $y$ of M87U1. In particular N86V1 applies to all its
required weighted averages.
% END M87U2

% BEGIN M87U3
\paragraph{M87U3. Recovering the original pointwise dimensions from the model.}
Suppose the prescribed model down map is injective in the following
sense: on every good simultaneous eigenvector with the distinguished
string its kernel is zero. For $f\in E_x$ with native downward
lift zero, apply L86M1 to obtain
$\mathrm{thG}(Af)=B\theta f=0$. If $f\ne0$, M87U1--2 put $Af$
in the class of nonzero vectors to which this injectivity applies.
This is impossible. Therefore the native down operator $T_i^*$
has zero kernel on $E_x$ at every good point, including those with
upper multiplicity greater than one.

Write $A_x$ for the native lower parameter subspace in F86D1.
One already has $T_i^*E_x\subseteq A_x$ by the lower spanning
theorem F86I2 and V86E2. Thus the injectivity just obtained
gives $\dim E_x\le\dim A_x$. Conversely $T_i$ is injective
and takes $A_x$ into $E_x$. Hence $T_iA_x=E_x$ individually
for every good point. Summing recovers equality in F86D1.
Empty good support has zero sum on both sides. F86D2 now gives
both original fields $\mathrm{what}_2$ and $\mathrm z$.
The proof needs injectivity only on these original vectors,
so it is equally valid if the stated model assertion is
established directly for the vectors of M87U1--2.
% END M87U3
